% Somas parciais S_N = \sum_{n=1}^N 1/n^2 (valores calculados e conferidos;
% 5 casas decimais) e o limite \pi^2/6 \approx 1,64493.
\begin{tikzpicture}
  \begin{axis}[
      width=0.82\linewidth, height=0.45\linewidth,
      xmin=0, xmax=21, ymin=0.9, ymax=1.76,
      xlabel={$N$}, ylabel={$S_N$},
      xtick={1,5,10,15,20},
      ytick={1,1.2,1.4,1.6},
      yticklabel style={/pgf/number format/.cd, use comma, fixed, precision=1},
      axis lines=left, tick align=outside,
      every axis plot/.append style={thick},
    ]
    \addplot[only marks, mark=*, mark size=1.6pt] coordinates {
      (1,1.00000) (2,1.25000) (3,1.36111) (4,1.42361) (5,1.46361)
      (6,1.49139) (7,1.51180) (8,1.52742) (9,1.53977) (10,1.54977)
      (11,1.55803) (12,1.56498) (13,1.57089) (14,1.57600) (15,1.58044)
      (16,1.58435) (17,1.58781) (18,1.59089) (19,1.59366) (20,1.59616)
    };
    \addplot[dashed, domain=0:21] {pi^2/6}
      node[pos=0.04, above right] {$\pi^2/6$};
  \end{axis}
\end{tikzpicture}
